\section{Recovering the two Wach heights from the filtered representation}

Fix a prime $p$, an original good layer $(K,s,L)$, its coefficient
DVR $O=\mathcal O_L$, a crystalline representation $V$ of $G_K$,
and its given full stable lattice $T$. Write $d=\dim_L V$,
$K_0$ for the maximal unramified subfield of $K$, and $i$ for an
embedding of $K_0$ in $L$ in the source's cyclic order. Put
$A_i=O[[X]]$ and $N_i=e_i N$ for the corresponding Wach component.
Use $Q=\Phi_{p^s}(1+X)$ and write $t=p^{-s}\iota_B(t_B)$ for the de Rham period used in the layer expansion.
Use the compatible Wach, functor and comparison records constructed
in L78W13 and chosen in F78D2.

% BEGIN W79H1
\paragraph{W79H1. Crystalline comparison after restricting to the good level.}
Put $K_s=K_0(\mu_{p^s})$. Goodness gives $K\subset K_s$.
The extension $K_s/K_0$ is totally ramified: the translated
cyclotomic polynomial $\Phi_{p^s}(1+X)$ is Eisenstein at $p$
over the unramified $K_0$. Thus the two fields $K$ and $K_s$
have the same maximal unramified subfield.

Put $B=B_{\rm cris}$, $S=L\otimes_{\mathbb Q_p}B$ and
$M=V\otimes_{\mathbb Q_p}B$, a free $S$-module of rank $d$.
On this very Wach lattice use the period map $b_T$ of L78W4--L78W5.
For a Wach basis $n_1,\ldots,n_d$, its period columns
$b_T(n_j)$ form an $S$-basis. Indeed L78W3 expresses
$x^a(1\otimes t_i)$, for each $O$-basis vector $t_i$ of $T$, as
an integral linear combination over $A_0$ of the $j(n_j)$,
with $x=1\otimes\pi_1$. Applying the cleared period identity
of L78W5 and dividing by the unit $\beta(x)^a$ gives $t_i\otimes1$
in the $S$-span of the $b_T(n_j)$. Explicitly one may first
multiply this relation by a common power of $x$ placing every
$j(n_j)$ in the image of the integral tensor module, where the
period map is defined. A square matrix spanning $S^d$ has a
right inverse and hence has unit determinant.

Multiply this row of columns by the normalized matrix $J(U)$ of
L78W9--L78W11. Its inverse is $J(U^{-1})$ because $J$ extends
through inversion of the logarithm, whose crystalline value is a
unit. The resulting ordered columns $d_1,\ldots,d_d$ are an
$S$-basis fixed by $G_K$: the action on the basis row has matrix
$J(G_g)$ and $G_g g(U)=U$. This construction needs the given
lattice and period coefficient map. It is independent of the
expansion maps used below.

Compute fixed vectors in this basis. For $G_K$ the fixed
scalars in $B$ are exactly $K_0$, by the $\mathrm{invB}$ field
of the given period datum. On $G_{K_s}$ the period action is the
restriction of the same absolute action. Its $\mathrm{invB}$ field
at $K_s$ gives the same $K_0$ by total ramification above.
Coefficientwise in a $\mathbb Q_p$-basis of $L$ this also proves
$S^{G_K}=S^{G_{K_s}}=L\otimes_{\mathbb Q_p} K_0$.
The invariant modules in $M$ for both groups are therefore
free of rank $d$ over this algebra, in the displayed basis. This
gives $\mathrm{IsCrysOn}$ with the original restricted action.
Since $L$ splits $K_0$, the algebra is $\prod_i L$, and in
particular the crystalline space $D_i$ has dimension $d$ for
every $i$. Extending the same basis to
$L\otimes_{\mathbb Q_p}B_{\rm dR}$ and using
$B_{\rm dR}^{G_K}=K$ gives
$D_{\rm dR}=K\otimes_{K_0}D_{\rm cris}$. Goodness splits $K$ in
$L$, so for each $\tau:K\hookrightarrow L$ above $i$, the
de Rham $\tau$-component receives $D_i$ isomorphically.
% END W79H1

% BEGIN W79H2
\paragraph{W79H2. A primitive cyclotomic expansion detects the filtered lattice.}
Fix a primitive $p^s$-th root $\zeta\in L$ and a component $i$.
There is a compatible embedding $\iota:L\hookrightarrow B_{\rm dR}$
which sends $\zeta$ to the chosen period root, takes the $i$th
copy of $K_0$ to its period inclusion, and determines
$\tau:K\hookrightarrow L$. Indeed the unramified field $K_0$
and the totally ramified cyclotomic field are linearly disjoint.
Map the subfield $\mathrm{emb}_i(K_0)(\zeta)$ of $L$ onto $K_s$
by the inverse of $\mathrm{emb}_i$ on $K_0$ and the prescribed
map on $\zeta$. The inverse of this isomorphism restricts to
$\tau:K\to L$, since $K\subset K_s$. Extend the isomorphism to
an embedding of $L$ into the algebraic closure of $\mathbb Q_p$,
then compose it with the tied algebraic period inclusion. This
gives the stated $\iota$, with $\iota\circ\tau$ the prescribed
inclusion of $K$.

Use the parameter $Y$ for the expansion $1+X=\zeta\exp(Y)$.
Let $\Lambda_{i,\zeta}\subset D_i((Y))$ be the $L[[Y]]$-span
of the Laurent expansions of $N$, using L78W11--L78W12.
The expansions of any one Wach basis generate this span, since
expansion takes multiplication by an integral series to multiplication
by its Taylor series. All generators have a finite pole bound.
After $Y\mapsto t$, scalar extension of $\Lambda_{i,\zeta}$ to
$B_{\rm dR}^+$ in the $\iota$-component is exactly
$V\otimes_{L,\iota}B_{\rm dR}^+$, under the period comparison.
To see this, use the two inclusions in L78W3, whose common clearing
factor is a power of $\pi_1$. Its de Rham image has nonzero residue
$\zeta_p-1$. Both this image and its inverse are therefore in
$B_{\rm dR}^+$. Map an integral lift of the cleared basis vector
through the de Rham period map and divide by this unit power.
L78W12 identifies the result in the coefficient component with
the evaluated Laurent expansion: after clearing its pole the
difference from the partial series of length $M$ lies in
$\mathrm{Fil}^{M-c}$ for a fixed $c$. Let $M$ tend to infinity.
Completeness and separatedness identify the two values.
The first inclusion puts every such column in the integral
period lattice. In the other direction the second inclusion
expresses the cleared vectors of $T$ as integral combinations
of those columns. The integral period coefficients map to
$B_{\rm dR}^+$. Dividing by the same unit there proves the spanning
assertion, also after projection to $\iota$.

Evaluation $L[[Y]]\to B_{\rm dR}^+$ here means
$\sum a_jY^j\mapsto\sum\iota(a_j)t^j$.
It preserves order since each nonzero algebraic coefficient is a
$t$-adic unit. The coordinate matrix of the $d$ columns relative
to a basis of $D_i$ consequently has nonzero determinant in
$L((Y))$, since its evaluated columns span in dimension $d$.
Thus they give a basis of a full Laurent lattice. For $u\in D_i((Y))$
solve for its coordinates in that basis over $L((Y))$.
The coordinates are in $L[[Y]]$ exactly when their evaluations
belong to $B_{\rm dR}^+$, by order preservation. In particular,
for a constant $d_0\in D_i$ and $k\in\mathbb Z$,
\[
 Y^k d_0\in\Lambda_{i,\zeta}
     \quad\Longleftrightarrow\quad
 d_0\in\mathrm{Fil}^{-k}D_{{\rm dR},\tau}.
\]
Here the filtration on the right is the one induced on the
$\tau$-component by the de Rham comparison.
Indeed it tests whether the period vector lies in
$t^{-k}(V\otimes B_{\rm dR}^+)$ at $\iota$. Testing one such
embedding suffices: the other embeddings inducing the same $\tau$
are conjugate under $G_K$, which fixes the de Rham vector and
preserves every filtration step. Splitting the finite separable
coefficient algebra over $B_{\rm dR}$ makes the component test
equivalent to the original filtration on that summand.
% END W79H2

% BEGIN W79H3
\paragraph{W79H3. Grading and the standard lattice below the critical level.}
For $\zeta^{p^s}=1$, the comparison with Frobenius sends the
expansion of a Frobenius-translated vector to the expansion
at $(i-1,\zeta^p)$ with coefficient in degree $k$ sent to
$p^k\varphi(d_k)$. This holds after multiplying by the expansion
of any power of $Q$ used to put that vector back in $N$.
Frobenius maps $D_{i-1}$ bijectively onto $D_i$.
Indeed it sends the idempotent for $i-1$ to the one for $i$;
in the normalized basis of W79H1 its coefficient matrix on that
component is the invertible constant $C_i$ of L78W9--L78W11. The substitution $Y\mapsto pY$
is an automorphism of $L[[Y]]$. Thus, over $L((Y))$, the two
expansion matrices on successive levels have the same rank:
the original two Wach inclusions identify $N_i$ and the span of
$\varphi N_{i-1}$ after inverting $Q$, and the expansion of $Q$
is nonzero. Every lower root is the $p$-th power, iterated a
suitable number of times, of a primitive root in $L$.
W79H2 therefore gives full rank at every lower root. At a root
of order dividing $p^{s-1}$ the expansion of $Q$ is a unit,
so the two lattices themselves agree under this Frobenius transport.

Each of these full lattices $\Lambda$ lies between
$Y^bD_i[[Y]]$ and $Y^aD_i[[Y]]$ for integers $a\le b$,
by bounding the orders of an expansion matrix and its inverse.
It is stable under $Y\mapsto uY$ for the cyclotomic generator
$u$ fixing the $p^s$-th roots. Its powers $u^j$ are distinct
for integer $j$, since $v_p(u-1)=s$ with $s\ge2$ for $p=2$.
For odd $p$, $s\ge1$; in both cases the $p$-adic logarithm is
injective on this subgroup. Stability follows from the Galois
coefficient formula and its inverse. In the finite quotient
$Y^aD_i[[Y]]/Y^bD_i[[Y]]$, Lagrange interpolation in this operator
projects onto each degree. Therefore $f\in\Lambda$ if and only if
each of its monomials belongs to $\Lambda$, for Laurent series
bounded below. For the reverse implication one takes the finite
part up to degree $b-1$ and then the entire tail in the lower
lattice $Y^bD_i[[Y]]$.

At $\zeta=1$ the same generator acts trivially on
$\Lambda/Y\Lambda$. On a Wach basis this follows from
$(\gamma-1)N\subset XN$ and $X=\exp Y-1\in YL[[Y]]$.
On the coefficient ring the substitution acts trivially modulo $Y$.
Put $D_k=\{d:Y^kd\in\Lambda\}$. These are increasing, zero for
$k<a$, and equal to $D_i$ for $k\ge b$. The quotient
$D_k/D_{k-1}$ contributes to $\Lambda/Y\Lambda$ with eigenvalue
$u^k$. Hence it vanishes except in degree zero. Grading now gives
$\Lambda_{i,1}=D_i[[Y]]$. Propagate this equality upward by
the unit Frobenius transports at roots of order $<p^s$. This
proves standardness at every lower root, including the predecessor
of each primitive root.
% END W79H3

% BEGIN W79H4
\paragraph{W79H4. Descent of the two integral height inequalities.}
Let all jumps at the embeddings $\tau$ above $i$ be contained in
$[-M_i,h_i]$. At a primitive root W79H2 and W79H3 give
\[
  Y^{M_i}D_i[[Y]]\subseteq\Lambda_{i,\zeta}
                 \subseteq Y^{-h_i}D_i[[Y]].
\]
For the left inclusion the filtration on constants is full in
degree $-M_i$: it is exhaustive and has no earlier jump. For
the right inclusion it is zero in degree $h_i+1$ by separatedness
and absence of later jumps, so grading excludes every term of
degree below $-h_i$. At the same root the expansion of the
lattice $P_i$ generated by $\varphi N_{i-1}$ is $D_i[[Y]]$,
by W79H3 applied at $\zeta^p$ and Frobenius transport.
The expansion of $Q$ equals $Y$ times a unit, since the roots
of the cyclotomic polynomial are simple and the coordinate
$\zeta\exp Y-1$ has nonzero linear coefficient.

It remains to recover integral, as opposed to rational Taylor,
inclusions. The two lattices $N_i,P_i$ are free of rank $d$ with
the indicated bases, and are equal on inverting $Q$ by the
initial, unquantified Wach height conditions. For $P_i$, its
generating columns are independent after inverting $Q$ because
their square matrix spans there; this gives freeness before
localizing as well. For an entry of a change of basis in either
direction write it as $f/Q^a$, $f\in O[[X]]$. At every primitive
root its Taylor expansion has the asserted order. Concretely,
if a numerator $g\in O[[X]]$ has Taylor order at least $a$ at
all these roots, $Q^a$ divides $g$ in $O[[X]]$. Indeed each
$z=\zeta-1$ belongs to the maximal ideal of $O$. Division of an
integral convergent-at-$z$ power series by $X-z$ has remainder
$g(z)$ and integral quotient: the coefficient of degree $n$ of
the quotient is $\sum_{j>n}g_j z^{j-n-1}$.
The remainder vanishes at each required step. Repeating at one
root and then the others preserves the orders at distinct roots
since the earlier linear factors have nonzero value there.
Their product is the monic $Q^a$. Apply this test to the
coordinates of $Q^{M_i}P_i$ in $N_i$ and of $Q^{h_i}N_i$ in
$P_i$ after clearing their denominators. The Taylor inclusions
give the required orders in the two bases. This proves
\[
 Q^{M_i}P_i\subseteq N_i,\qquad Q^{h_i}N_i\subseteq P_i
\]
on the original integral module. In particular the argument
does not require differences of distinct cyclotomic roots to
be units of $O$.
% END W79H4

% BEGIN W79H5
\paragraph{W79H5. The complete local Wach hypothesis on the stored module.}
Suppose now $T$ is crystalline and every jump in the full de Rham
jump set lies in $[-M_0,h_0]$. An embedding-specific jump is in
that full set: intersect equality of the two global filtration
steps with the embedding component. Use W79H4 with the two
constant bounds and put $N=\bigoplus_i N_i$. The first inclusion
says $Q^{M_0}\varphi N\subseteq N$ and the second gives
$Q^{h_0}N\subseteq A\varphi(N)$, where $A=\prod_i A_i$.

Take the basis of L78W1. Its images under the original map $j$
are a basis over the Witt comparison ring by L78W3: they generate
the comparison module by the two scalar extensions, and a square
surjective matrix over a commutative ring has unit determinant.
The action of the group at level $s$ preserves $N$, its difference
from the identity has image in $XN$, and lattice endomorphisms
preserve $N$. These are the action, difference and functoriality
laws of the chosen Wach record. Together with the basis and the
two displayed inclusions they give all six fields of
$\mathrm{WClauses}$. W79H1 supplies crystallinity on the specified
subgroup. Joining the original jump interval gives
$\mathrm{HypW}(T,M_0,h_0)$ on the recorded functor and comparison.
The prime, good layer and crystalline lattice were arbitrary in
this chosen family of compatible records. Thus their original
$\mathrm{WAll}(D,H,F)$ holds for the very $H,F$ and comparison
from F78D2. In particular it applies to the self coefficient
bridge and lattice of S79P1--S79P4.
% END W79H5
